\sectionkicker{Source-backed technical notes}
\subsection*{Inference \& Serving — cache・MoE・kernel・runtimeの最適化: Technical Notes}
本欄は記事本文で圧縮した一次資料上の情報を、比較・再検証しやすい形で整理したものである。時系列、確認済みの主張、留意点、一次資料URLを掲載する。
\medskip
\subsection*{Theme at a glance}
\addcontentsline{toc}{subsection}{Theme at a glance}
\begin{center}
\footnotesize
\begin{tabularx}{\linewidth}{@{}p{0.20\linewidth}p{0.10\linewidth}p{0.14\linewidth}X@{}}
\toprule
資料 & 位置づけ & 種別 & 時系列 \\
\midrule
vllm-project/vllm May 2026 release series & 主要資料 & Framework & 2026-05-04 (Framework公開); 2026-05-15 (Framework公開); 2026-05-29 (Framework公開) \\
sgl-project/sglang May 2026 release series & 主要資料 & Framework & 2026-05-05 (Framework公開); 2026-05-16 (Framework公開) \\
Release v0.6.12 & 主要資料 & Framework & 2026-05-29 (Framework公開) \\
Adaptive KV Cache Reuse for Fast Long-Context LLM Serving & 補足資料 & 論文 & 2026-05-20 (論文公開) \\
ReMoE: Boosting Expert Reuse through Router Fine-Tuning in Memory-Constrained MoE LLM Inference & 補足資料 & 論文 & 2026-05-26 (論文公開) \\
\bottomrule
\end{tabularx}
\end{center}
\normalsize

\begin{technicalnote}{vllm-project/vllm May 2026 release series}{主要資料}
\begingroup
% reader-facing Technical Notes break policy
\widowpenalty=10000
\clubpenalty=10000
\displaywidowpenalty=10000
\begin{tabularx}{\linewidth}{@{}>{\bfseries}p{0.22\linewidth}X@{}}
組織 & vllm-project/vllm \\
種別 & Framework \\
時系列 & 2026-05-04 (Framework公開); 2026-05-15 (Framework公開); 2026-05-29 (Framework公開) \\
\end{tabularx}
\smallskip
{\bfseries 一次資料から整理したtechnical points}
\begin{itemize}[leftmargin=1.5em,itemsep=0.35em]
\item \textbf{一次情報で確認できる事実}: 保存済み一次資料では「v0.20.1」が2026-05-04に確認できる。
\item \textbf{一次情報で確認できる事実}: 保存済み一次資料では「v0.21.0」が2026-05-15に確認できる。
\item \textbf{一次情報で確認できる事実}: 保存済み一次資料では「v0.22.0」が2026-05-29に確認できる。
\end{itemize}
\begin{minipage}{\linewidth}
% reader-facing Technical Notes coherent tail group
\begin{itemize}[leftmargin=1.5em,itemsep=0.35em]
\item \textbf{Project claim}: vLLMのMay release series（v0.20.1、v0.21.0、v0.22.0）は、DeepSeek V4 hardening、KV/offload、attention backend、Model Runner V2、speculative/disaggregated servingなどを扱っている。
\end{itemize}
{\bfseries 読む際の境界}
\begin{itemize}[leftmargin=1.5em,itemsep=0.35em]
\item \textbf{分析上の留意点}: 公式repositoryのrelease notesは公開時系列と記載された実装変更を確認する一次資料だが、projectのperformance claimやhardwareをまたぐ一般化は本号で独立再現していない。
\end{itemize}
\begin{samepage}
{\bfseries 一次資料}
\begingroup\sloppy
\begin{itemize}[leftmargin=1.5em,itemsep=0.25em]
\item {\scriptsize\url{https://github.com/vllm-project/vllm/releases/tag/v0.20.1}}
\item {\scriptsize\url{https://github.com/vllm-project/vllm/releases/tag/v0.21.0}}
\item {\scriptsize\url{https://github.com/vllm-project/vllm/releases/tag/v0.22.0}}
\end{itemize}
\endgroup
\end{samepage}
\end{minipage}
\endgroup
\end{technicalnote}
\medskip

\begin{technicalnote}{sgl-project/sglang May 2026 release series}{主要資料}
\begingroup
% reader-facing Technical Notes break policy
\widowpenalty=10000
\clubpenalty=10000
\displaywidowpenalty=10000
\begin{tabularx}{\linewidth}{@{}>{\bfseries}p{0.22\linewidth}X@{}}
組織 & sgl-project/sglang \\
種別 & Framework \\
時系列 & 2026-05-05 (Framework公開); 2026-05-16 (Framework公開) \\
\end{tabularx}
\smallskip
{\bfseries 一次資料から整理したtechnical points}
\begin{itemize}[leftmargin=1.5em,itemsep=0.35em]
\item \textbf{一次情報で確認できる事実}: 保存済み一次資料では「v0.5.11」が2026-05-05に確認できる。
\item \textbf{一次情報で確認できる事実}: 保存済み一次資料では「v0.5.12」が2026-05-16に確認できる。
\end{itemize}
\begin{minipage}{\linewidth}
% reader-facing Technical Notes coherent tail group
\begin{itemize}[leftmargin=1.5em,itemsep=0.35em]
\item \textbf{Project claim}: SGLangのMay releases v0.5.11 / v0.5.12は、speculative decoding、disaggregated prefix caching、model support、DeepSeek V4 inference pathなどを扱っている。
\end{itemize}
{\bfseries 読む際の境界}
\begin{itemize}[leftmargin=1.5em,itemsep=0.35em]
\item \textbf{分析上の留意点}: 公式repositoryのrelease notesは公開時系列と記載された実装変更を確認する一次資料だが、projectのperformance claimやhardwareをまたぐ一般化は本号で独立再現していない。
\end{itemize}
\begin{samepage}
{\bfseries 一次資料}
\begingroup\sloppy
\begin{itemize}[leftmargin=1.5em,itemsep=0.25em]
\item {\scriptsize\url{https://github.com/sgl-project/sglang/releases/tag/v0.5.11}}
\item {\scriptsize\url{https://github.com/sgl-project/sglang/releases/tag/v0.5.12}}
\end{itemize}
\endgroup
\end{samepage}
\end{minipage}
\endgroup
\end{technicalnote}
\medskip

\begin{technicalnote}{Release v0.6.12}{主要資料}
\begingroup
% reader-facing Technical Notes break policy
\widowpenalty=10000
\clubpenalty=10000
\displaywidowpenalty=10000
\begin{tabularx}{\linewidth}{@{}>{\bfseries}p{0.22\linewidth}X@{}}
組織 & flashinfer-ai/flashinfer \\
種別 & Framework \\
時系列 & 2026-05-29 (Framework公開) \\
\end{tabularx}
\smallskip
{\bfseries 一次資料から整理したtechnical points}
\begin{itemize}[leftmargin=1.5em,itemsep=0.35em]
\item \textbf{一次情報で確認できる事実}: 保存済み一次資料では「Release v0.6.12」が2026-05-29に確認できる。
\end{itemize}
\begin{minipage}{\linewidth}
% reader-facing Technical Notes coherent tail group
\begin{itemize}[leftmargin=1.5em,itemsep=0.35em]
\item \textbf{Project claim}: FlashInfer v0.6.12は、NVFP4 quantization、MoE kernel、GQA/MLA path、model/backend supportなど、low-level servingに関わるinference-kernel更新を含む。
\end{itemize}
{\bfseries 読む際の境界}
\begin{itemize}[leftmargin=1.5em,itemsep=0.35em]
\item \textbf{分析上の留意点}: 公式repositoryのrelease notesは公開時系列と記載された実装変更を確認する一次資料だが、projectのperformance claimやhardwareをまたぐ一般化は本号で独立再現していない。
\end{itemize}
\begin{samepage}
{\bfseries 一次資料}
\begingroup\sloppy
\begin{itemize}[leftmargin=1.5em,itemsep=0.25em]
\item {\scriptsize\url{https://github.com/flashinfer-ai/flashinfer/releases/tag/v0.6.12}}
\end{itemize}
\endgroup
\end{samepage}
\end{minipage}
\endgroup
\end{technicalnote}
\medskip

\begin{technicalnote}{Adaptive KV Cache Reuse for Fast Long-Context LLM Serving}{補足資料}
\begingroup
% reader-facing Technical Notes break policy
\widowpenalty=10000
\clubpenalty=10000
\displaywidowpenalty=10000
\begin{tabularx}{\linewidth}{@{}>{\bfseries}p{0.22\linewidth}X@{}}
組織 & - \\
種別 & 論文 \\
時系列 & 2026-05-20 (論文公開) \\
\end{tabularx}
\smallskip
{\bfseries 一次資料から整理したtechnical points}
\begin{itemize}[leftmargin=1.5em,itemsep=0.35em]
\item \textbf{一次情報で確認できる事実}: 保存済み一次資料では「Adaptive KV Cache Reuse for Fast Long-Context LLM Serving」が2026-05-20に確認できる。
\end{itemize}
\begin{minipage}{\linewidth}
% reader-facing Technical Notes coherent tail group
\begin{itemize}[leftmargin=1.5em,itemsep=0.35em]
\item \textbf{Author claim}: CacheTuneは、strict prefix一致を越えてKV cacheを再利用し、long-context servingでattention関係を保つため必要なcontextを適応的に再計算する方式を提案している。
\end{itemize}
{\bfseries 読む際の境界}
\begin{itemize}[leftmargin=1.5em,itemsep=0.35em]
\item \textbf{分析上の留意点}: arXivの原メタデータ／abstractは一次論文資料として扱うが、実験結果と一般化可能性はAuthor claimであり、本号では独立再現していない。
\end{itemize}
\begin{samepage}
{\bfseries 一次資料}
\begingroup\sloppy
\begin{itemize}[leftmargin=1.5em,itemsep=0.25em]
\item {\scriptsize\url{http://arxiv.org/abs/2605.24022v1}}
\end{itemize}
\endgroup
\end{samepage}
\end{minipage}
\endgroup
\end{technicalnote}
\medskip

\begin{technicalnote}{ReMoE: Boosting Expert Reuse through Router Fine-Tuning in Memory-Constrained MoE LLM Inference}{補足資料}
\begingroup
% reader-facing Technical Notes break policy
\widowpenalty=10000
\clubpenalty=10000
\displaywidowpenalty=10000
\begin{tabularx}{\linewidth}{@{}>{\bfseries}p{0.22\linewidth}X@{}}
組織 & - \\
種別 & 論文 \\
時系列 & 2026-05-26 (論文公開) \\
\end{tabularx}
\smallskip
{\bfseries 一次資料から整理したtechnical points}
\begin{itemize}[leftmargin=1.5em,itemsep=0.35em]
\item \textbf{一次情報で確認できる事実}: 保存済み一次資料では「ReMoE: Boosting Expert Reuse through Router Fine-Tuning in Memory-Constrained MoE LLM Inference」が2026-05-26に確認できる。
\end{itemize}
\begin{minipage}{\linewidth}
% reader-facing Technical Notes coherent tail group
\begin{itemize}[leftmargin=1.5em,itemsep=0.35em]
\item \textbf{Author claim}: ReMoEは、memory-constrained inferenceで一部expertのみをcacheできる条件下でexpert reuseを増やすため、MoE routingをfine-tuningする方式を提案している。
\end{itemize}
{\bfseries 読む際の境界}
\begin{itemize}[leftmargin=1.5em,itemsep=0.35em]
\item \textbf{分析上の留意点}: arXivの原メタデータ／abstractは一次論文資料として扱うが、実験結果と一般化可能性はAuthor claimであり、本号では独立再現していない。
\end{itemize}
\begin{samepage}
{\bfseries 一次資料}
\begingroup\sloppy
\begin{itemize}[leftmargin=1.5em,itemsep=0.25em]
\item {\scriptsize\url{http://arxiv.org/abs/2605.27081v1}}
\end{itemize}
\endgroup
\end{samepage}
\end{minipage}
\endgroup
\end{technicalnote}
\medskip
